\begin{textsample}{POS Dim 4 – pa – Score 12.00 – pa/pa\_inf\_8.txt}  \label{ex:f4_pos_015}
3.1.2.5 Educação Infantil e a Criança Quilombola Falar na criança no contexto da Educação Escolar Quilombola é, antes de tudo, pensar em uma educação que possa valorizar o passado e reinventar o presente de uma \textbf{população} que teve seus direitos \textbf{negados} por longos anos.
Desse modo, é preciso que, desde a Educação Infantil, a história, a memória, a tecnologia, o \textbf{território} e os conhecimentos dos \textbf{quilombos} sejam reconhecidos e considerados no currículo escolar, conforme as Diretrizes Curriculares Nacionais para a Educação Escolar Quilombola ( BRASIL, 2012b ).
A história dos \textbf{quilombos} é uma história de luta e de resistência, sobretudo no tocante à língua.
A palavra \textbf{quilombo}, derivada do banto, pode estar relacionada à aldeia, acampamento ou similares.
O banto se refere ao grupo africano étnico-linguístico, composto de várias línguas chamadas línguas bantas que se caracterizam pela utilização de prefixos.
A manutenção das línguas dos povos escravizados era uma forma de resistência, haja vista que os donos dos mesmos não conheciam determinadas línguas, era também uma forma de preservação da \textbf{identidade} desses povos.
O Conselho Ultramarino de 1740 definiu \textbf{quilombo} como “ toda habitação de negros fugidos, que passam de cinco, em parte despovoada, ainda que não tenham ranchos levantados e nem se achem pilões neles ” ( MOURA, 1997, p. 87 ).
Os \textbf{quilombos} representavam um núcleo de resistência, que surgiam em qualquer espaço onde houvesse escravismo, dos séculos XVI ao XIX no Brasil, daí surge a expressão “ quilombolas ” para classificar os negros pertencentes aos \textbf{quilombos}.
As comunidades quilombolas se constituem de forma coletiva, ou seja, a titulação da terra se dá em nome de uma associação seguido de uma lista com o nome de todos que ali residem, assim a \textbf{identidade} quilombola está ligada também à ideia de pertencimento que se estabelece muito além do que laços de sangue e se fortalecem por meio dos valores, costumes e experiência de \textbf{discriminação}, compartilhados por um grupo com um sentimento em comum.
A titulação das terras está para além da regularização fundiária, pois formaliza uma política de cidadania que mantém vivos os valores e a história de um povo.
Vale ressaltar que, o \textbf{estado} do Pará foi a primeira unidade da federação a titular terras em favor das comunidades remanescentes de \textbf{quilombos} no ano de 1995, em Oriximiná.
Além disso, ocupa a colocação de \textbf{estado} que mais demarcou \textbf{territórios} quilombolas no Brasil ( MARQUES ; MALCHER, 2009 ).
Segundo a Comissão Pró-Índio de São Paulo, organização não governamental, o \textbf{estado} do Pará titulou 63 terras quilombolas das 85 existentes no \textbf{Estado}.
Estudos recentes apontam para a existência de mais de quatrocentas comunidades.
A partir da divisão adotada pelo Instituto Brasileiro de Geografia e Estatística – IBGE –, que divide o \textbf{Estado} em seis mesorregiões, o Programa identificou a existência dessas comunidades em quatro delas : Baixo Amazonas, Marajó, \textbf{Nordeste} e Metropolitana de Belém ( MARQUES ; MALCHER, 2009 ).
O Programa Raízes, criado pelo Decreto de 11 de maio de 2000, no Pará, identificou 240 comunidades nas terras quilombolas, contudo os estudos atuais apontam para 400 comunidades, como é possível constatar acima.
No que se refere à educação quilombola, pode -se dizer que passou por um longo período de esquecimento diluída nas políticas da Educação Rural, sem nenhuma política pública e ou pedagógica que considerasse a sua especificidade.
No entanto, o resultado das mobilizações, tecidas no \textbf{bojo} dos movimentos sociais com destaque para o Movimento Negro e para o Movimento Quilombola, fez com que fosse delineado um movimento de discussões sobre mudanças no modelo de ensino para as escolas das comunidades quilombolas atendendo de forma específica e diferenciadas as crianças.
As muitas lutas tecidas pelos movimentos sociais culminaram na promulgação da Resolução no 08 de 20 de novembro de 2012 que institui as Diretrizes Curriculares Nacionais para Educação Escolar Quilombola ; a publicação dessa legislação pode ser considerada um dos marcos na luta do Movimento Negro e do Movimento Quilombola, pois ela consolida a Educação Escolar Quilombola como uma modalidade de ensino da Educação Básica.
No cumprimento da Educação Infantil como um dos níveis da Educação Básica, a Educação Escolar Quilombola deve ser desenvolvida de acordo com a Resolução CNE/CEB No 4/2010, que define as Diretrizes Curriculares Nacionais para a Educação Básica e a Resolução CNE/CEB no 5/2009, que definiu as Diretrizes Curriculares para a Educação Infantil, bem como considerar os aspectos específicos dessas \textbf{populações} na vivência de suas infâncias, destacadas estas diretrizes e construída em conjunto com as comunidades a que pertencem.
A Educação Escolar Quilombola é desenvolvida em unidades educacionais inscritas em suas terras e culturas, requerendo pedagogia própria em respeito às especificidades étnico-cultural de cada comunidade e formação docente, observados os princípios constitucionais, a base nacional comum e os princípios que orientam a Educação básica brasileira.
Na estruturação e no funcionamento das escolas quilombolas, deve ser reconhecida e valorizada sua diversidade cultural ( BRASIL, 2012b, p. 42 ).
Nesse sentido, a Educação Infantil Quilombola deve estar pautada nas DCNEI ( BRASIL, 2010a ) em que são recomendados o cuidar e o educar como elementos indissociáveis nessa etapa da educação básica, além de serem um direito das crianças quilombolas e, portanto obrigatória a oferta pelo poder público para as crianças em idade pré-escolar de 4 e 5 anos.
Quanto às crianças em idade de creche, 0 a 03 anos, fica a critério das famílias decidirem pela matrícula, uma vez que é também o direito da criança a permanência junto ao seu grupo familiar e ao grupo comunitário de referência.
No que concerne à organização curricular da Educação Infantil Quilombola esta deve se dar de forma democrática e horizontal, visto que todos devem ter direito a voz e escuta, rompendo com a tradição de silêncio imposta a esse povo.
Logo, o atendimento educacional das crianças que vivem nessas comunidades precisa estar pautado nos saberes ora pertencentes a esse povo.
O conhecimento tradicional de cada comunidade deverá ser expresso a partir da participação das famílias e dos anciãos que são os especialistas nas tradições do seu povo.
Além disso, os saberes das comunidades remanescentes que fazem parte da história de cada um e que devem refletir a realidade e o contexto local de cada comunidade, devem permear o currículo escolar, visando a promover uma educação transformadora com práticas educacionais que assegurem a diversidade étnico \textbf{racial} da \textbf{população} ali existente.
A Educação Infantil Quilombola deve valorizar a história e a tradição do seu povo como elementos indispensáveis à formação para a cidadania e afirmação da \textbf{identidade} cultural, esta manifestada pelas crianças, sobretudo, por intermédio de seus modos de vida, das brincadeiras, do trabalho e das relações com os adultos e idosos.
As crianças quilombolas apresentam modos próprios de se relacionar com o cotidiano e conviver com a natureza e seus diferentes espaços.
As brincadeiras representam suas produções em nível cognitivo, afetivo e social, tanto dentro quanto fora da escola.
Nesse momento, constroem seu universo próprio a partir das suas interações com o contexto local ( POJO ; BARRETO, 2013 ).
Desse modo, os currículos da Educação Infantil Quilombola devem ser construídos levando em consideração a forma organizativa das comunidades, suas contribuições sociais, culturais, econômicas, políticas, e não somente suas vestimentas, rituais festivos, entre outros fatores que fazem parte do cotidiano dos mesmos.
Assim como as brincadeiras e a confecção de materiais pedagógicos próprios da cultura quilombola, devem fazer parte do currículo escolar, respeitando as particularidades locais, como tempo, espaços pedagógicos, condições climáticas, rotinas que propiciam a valorização da \textbf{identidade} quilombola, que os identifiquem, que permitam se sentir partícipes de sua história, que provoque a criança a se reconhecer como sujeito, conhecer suas origens, entender que a história do povo quilombola foi construída por meio de lutas, contra o \textbf{racismo}, as desigualdades sociais, a conquista da \textbf{territorialidade}, a saúde, a moradia, da educação ; enfim, sobretudo, a busca pela garantia de direitos.

% matched lemmas: bojo, discriminação, estado, identidade, negar, nordeste, população, quilombo, racial, racismo, territorialidade, território
\end{textsample}
